% Identical five-page content for all typography variants, v1.0.10 review.
% Reuses stable authoring commands; no geometry, color or style redefinitions.
\setcounter{chapter}{1}
\setcounter{section}{0}
\setcounter{equation}{0}
\setcounter{figure}{0}
\setcounter{workedexample}{0}
\setcounter{page}{1}
\BookRememberRunningChapter{Κίνηση σε μία διάσταση}
\markboth{Κίνηση σε μία διάσταση}{Κίνηση σε μία διάσταση}

% PAGE 1 — Physics theory and inline math.
\BookSection{Θεωρία}
\section{Περιγραφή της κίνησης}

% Source: chapters/demo-physics.tex:19-19
Για την περιγραφή της ευθύγραμμης κίνησης ορίζουμε έναν άξονα αναφοράς και θετική φορά. Η θέση ενός σώματος προσδιορίζεται από την αλγεβρική τιμή της συντεταγμένης του.

% Source: chapters/starter-chapter.tex:12-17
Η θέση ενός σώματος σε έναν άξονα δίνεται από τη συντεταγμένη του.
Για κίνηση από τη θέση $x_A$ στη θέση $x_B$, η μετατόπιση είναι
\begin{equation}\label{eq:starter-displacement}
  \Delta x=x_B-x_A.
\end{equation}


% Source: chapters/demo-assessment.tex:25-27
Το μέτρο της μέσης διανυσματικής ταχύτητας ενός σώματος
υπολογίζεται από τη σχέση $|\bar{\upsilon}|=\frac{|\Delta x|}{\Delta t}$.
Να αναφέρετε τις μονάδες SI των μεγεθών $\Delta x$ και $\Delta t$.

% Source: chapters/demo-physics.tex:24-26
\begin{KeyFormula}
\(\Delta x=x_{\mathrm{τελ}}-x_{\mathrm{αρχ}}\qquad\text{και}\qquad \upsilon=\dfrac{\Delta x}{\Delta t}\)
\end{KeyFormula}

\BookMinorHeading{Από το διάγραμμα στη φυσική ερμηνεία}

% Source: chapters/demo-physics.tex:42-42
Η μορφή της γραφικής παράστασης παρέχει πληροφορίες για την κινητική κατάσταση. Στο ακόλουθο δοκιμαστικό διάγραμμα η θέση αυξάνεται γραμμικά με τον χρόνο.

% Source: stress-content.tex:5-9
Η περιγραφή ενός φαινομένου αρχίζει με τον ορισμό των μεγεθών και των μονάδων
τους. Καταγράφουμε τα δεδομένα, επιλέγουμε σύστημα αναφοράς και διατυπώνουμε
τις παραδοχές. Στη συνέχεια ελέγχουμε τη μαθηματική σχέση, το πρόσημο και τη
φυσική σημασία του αποτελέσματος. Η τελική απάντηση περιλαμβάνει κατάλληλη
μονάδα και σύντομη αιτιολόγηση.\par

% Source: stress-content.tex:13-14
\[x(t)=x_0+\upsilon_0t+\frac12at^2,\qquad
  \upsilon(t)=\upsilon_0+at.\]


\clearpage
% PAGE 2 — Both worked examples, including their original callouts.
\BookSection{Λυμένα παραδείγματα}

% Source: chapters/demo-physics.tex:115-143
\begin{workedexample}{Υπολογισμός μετατόπισης}
\begin{ProblemBox}
Ένα σώμα κινείται πάνω σε ευθεία από τη θέση \(x_1=-3\,\mathrm{m}\) στη θέση \(x_2=5\,\mathrm{m}\).\par
\begin{enumerate}[label=\BookGreek*),leftmargin=7mm,topsep=1pt,itemsep=0pt]
\item Να υπολογιστεί η μετατόπιση.
\item Να προσδιοριστεί η φορά της.
\end{enumerate}
\end{ProblemBox}
\begin{ThinkBox}
Ξεκινώ από τον ορισμό της μετατόπισης και εξετάζω το πρόσημο του αποτελέσματος. Το θετικό πρόσημο δηλώνει κατεύθυνση προς τη θετική φορά του άξονα.
\end{ThinkBox}
\SolutionHeading
\BookListLabel{α)} Από τον ορισμό της μετατόπισης:
\[\Delta x=x_2-x_1=5-(-3)=8\,\mathrm{m}.\]
\BookListLabel{β)} Η μετατόπιση έχει θετική αλγεβρική τιμή, επομένως κατευθύνεται προς τη θετική φορά του άξονα.
\end{workedexample}
\begin{workedexample}{Σταθερή ταχύτητα}
\begin{ProblemBox}
Ένα κινητό διανύει μετατόπιση \(\Delta x=60\,\mathrm{m}\) σε χρονικό διάστημα \(\Delta t=12\,\mathrm{s}\). Να υπολογιστεί η αλγεβρική τιμή της σταθερής ταχύτητας.
\end{ProblemBox}
\begin{ThinkBox}
Στην ευθύγραμμη ομαλή κίνηση χρησιμοποιώ τον λόγο \(\Delta x/\Delta t\) και ελέγχω τις μονάδες του αποτελέσματος.
\end{ThinkBox}
\SolutionHeading
\[\upsilon=\frac{\Delta x}{\Delta t}=\frac{60\,\mathrm{m}}{12\,\mathrm{s}}=5\,\mathrm{m/s}.\]
\end{workedexample}


\clearpage
% PAGE 3 — Dense math plus explicitly marked missing-operator diagnostics.
\BookSection{Μαθηματικοί συμβολισμοί}
\section{Η συνάρτηση πρώτου βαθμού}

% Source: chapters/demo-math.tex:16-16
Ο τύπος \(f(x)=\alpha x+\beta\), με \(\alpha\neq0\), περιγράφει μια συνάρτηση πρώτου βαθμού. Η γραφική της παράσταση είναι ευθεία.

% Source: stress-content.tex:379-386
\begin{align}
 f(x)&=\alpha x+\beta,\label{eq:stress-function}\\
 f(x+h)-f(x)&=\alpha h,\\
 \lim_{h\to0}\frac{f(x+h)-f(x)}{h}&=\alpha,\\
 \int_a^b(\alpha x+\beta)\,\mathrm dx
 &=\frac{\alpha}{2}(b^2-a^2)+\beta(b-a).
\end{align}
Η σχέση~\eqref{eq:stress-function} συνδέει τα παραδείγματα του κεφαλαίου.

% Source: chapters/demo-math.tex:18-20
\begin{definitionbox}[Συνάρτηση πρώτου βαθμού]
\(f(x)=\alpha x+\beta\), όπου \(\alpha\neq0\) και \(\alpha,\beta\in\mathbb R\).
\end{definitionbox}


\BookMinorHeading{Συμπληρωματική τυπογραφική δοκιμή}
Οι επόμενες σχέσεις προστίθενται μόνο για τον έλεγχο ριζών, αθροισμάτων,
διαδοχικών δεικτών και ελληνικών χαρακτήρων.
\begin{align*}
 \sqrt{x^2+y^2}&=\sqrt{\frac{a^2+b^2}{1+t^2}}, &
 S_n&=\sum_{k=1}^{n}\frac{1}{k^2},\\
 x_{i_j}^{\,2}&=\frac{a_{i+1}-a_i}{\Delta t}, &
 \vec F&=m\vec a.
\end{align*}
Λατινικά: $a,\allowbreak b,\allowbreak c,\allowbreak x,\allowbreak y,\allowbreak z,\allowbreak A,\allowbreak B,\allowbreak F,\allowbreak N$.
Ελληνικά: $\alpha,\allowbreak \beta,\allowbreak \gamma,\allowbreak \delta,\allowbreak \theta,\allowbreak \lambda,\allowbreak \mu,\allowbreak \nu,\allowbreak \pi,\allowbreak \rho,\allowbreak \sigma,\allowbreak 
\tau,\allowbreak \phi,\allowbreak \varphi,\allowbreak \chi,\allowbreak \psi,\allowbreak \omega,\allowbreak \Gamma,\allowbreak \Delta,\allowbreak \Theta,\allowbreak \Sigma,\allowbreak \Omega$.

Δοκιμή χημικής γραφής: $\mathrm{H_2O}$, $\mathrm{H_2SO_4}$,
$\mathrm{Na^+}$, $\mathrm{SO_4^{2-}}$ και $10^{-3}\,\mathrm{mol\,L^{-1}}$.
Ζεύγη σύγκρισης: $\nu/v$, $\rho/p$, $\mu/u$ και $\alpha/a$.

\clearpage
% PAGE 4 — Standard callouts, including dense inline math in a small context.
\BookSection{Διδακτικά κουτιά}
\section{Επιλογή σχέσης και έλεγχος}

% Source: stress-content.tex:5-9
Η περιγραφή ενός φαινομένου αρχίζει με τον ορισμό των μεγεθών και των μονάδων
τους. Καταγράφουμε τα δεδομένα, επιλέγουμε σύστημα αναφοράς και διατυπώνουμε
τις παραδοχές. Στη συνέχεια ελέγχουμε τη μαθηματική σχέση, το πρόσημο και τη
φυσική σημασία του αποτελέσματος. Η τελική απάντηση περιλαμβάνει κατάλληλη
μονάδα και σύντομη αιτιολόγηση.\par

% Source: chapters/demo-physics.tex:38-40
\begin{AttentionBox}
Η μετατόπιση είναι διανυσματικό μέγεθος, ενώ το διάστημα είναι μονόμετρο. Δεν ταυτίζονται σε κάθε κίνηση.
\end{AttentionBox}

% Source: chapters/demo-physics.tex:138-140
\begin{ThinkBox}
Στην ευθύγραμμη ομαλή κίνηση χρησιμοποιώ τον λόγο \(\Delta x/\Delta t\) και ελέγχω τις μονάδες του αποτελέσματος.
\end{ThinkBox}

% Source: chapters/demo-math.tex:57-59
\begin{ThinkBox}
Αντικαθιστώ στη μεταβλητή \(x\) την τιμή που μου ζητείται, χωρίς να αλλάζω τον τύπο της συνάρτησης.
\end{ThinkBox}


\begin{ThinkBox}
Συμπληρωματική τυπογραφική δοκιμή μέσα στο κουτί:
για $r=\sqrt{x^2+y^2}$, $q=\frac{a_{i+1}-a_i}{\Delta t}$ και
$S_n=\sum_{k=1}^{n}\frac{1}{k^2}$, εξετάζουμε τους δείκτες, τα κλάσματα
και την απόσταση των διαδοχικών γραμμών. Οι $\alpha$, $\beta$ και $\upsilon$
παραμένουν ευδιάκριτες από τα λατινικά $a$, $b$ και $v$.
\end{ThinkBox}

\clearpage
% PAGE 5 — Actual stable TikZ geometry and math labels.
\BookSection{Διαγράμματα και σύμβολα}
\section{Θέση, μετατόπιση και χρόνος}

\begin{center}

% Source: chapters/starter-chapter.tex:27-40
  \begin{tikzpicture}[x=.8cm,y=.8cm]
    \draw[book/axis] (-3.2,0)--(3.5,0)
      node[book/physlabel,right] {$x\,(\mathrm{m})$};
    \foreach \x in {-2,0,2}{%
      \draw[book/tick] (\x,.08)--(\x,-.08);
      \node[book/ticklabel,below=2pt] at (\x,-.08) {$\x$};
    }
    \node[book/endpoint] at (-2,0) {};
    \node[book/endpoint] at (2,0) {};
    \node[book/physlabel,above] at (-2,0) {$A$};
    \node[book/physlabel,above] at (2,0) {$B$};
    \draw[book/mainvector] (-2,.65)--(2,.65)
      node[book/physlabel,midway,above] {$\Delta x$};
  \end{tikzpicture}

\end{center}

\captionof{figure}{Μετατόπιση κινητού από την αρχική θέση Α στην τελική θέση Β.}
\label{fig:displacement}\label{fig:starter-displacement}

% Source: chapters/demo-physics.tex:42-42
Η μορφή της γραφικής παράστασης παρέχει πληροφορίες για την κινητική κατάσταση. Στο ακόλουθο δοκιμαστικό διάγραμμα η θέση αυξάνεται γραμμικά με τον χρόνο.

% Source: chapters/demo-physics.tex:43-53
\begin{center}
\begin{tikzpicture}[x=.65cm,y=.40cm]
\draw[book/axis] (0,0)--(8.0,0) node[book/physlabel,right] {$t\,(\mathrm{s})$};
\draw[book/axis] (0,0)--(0,7.0) node[book/physlabel,above] {$x\,(\mathrm{m})$};
\foreach \x in {2,4,6}{\draw[book/tick] (\x,.10)--(\x,-.10) node[book/ticklabel,below]{$\x$};}
\foreach \y in {2,4,6}{\draw[book/tick] (.10,\y)--(-.10,\y) node[book/ticklabel,left]{$\y$};}
\draw[book/maincurve] (0,1)--(7.2,6.4);
\draw[book/guide] (4,0)--(4,4);
\node[book/pointMain] at (4,4) {};
\end{tikzpicture}
\end{center}

\captionof{figure}{Η θέση αυξάνεται γραμμικά με τον χρόνο.}

Όπως φαίνεται στο Σχήμα~\ref{fig:displacement}, η μετατόπιση παριστάνεται με
διάνυσμα από την αρχική προς την τελική θέση.
